% !TeX program = pdflatex
% papers/estrela.tex — a estrela e o buraco negro como os dois sinais do trial.
% Compila sozinho:  pdflatex estrela.tex
\documentclass[11pt,a4paper]{article}
\usepackage[utf8]{inputenc}\usepackage[T1]{fontenc}\usepackage[portuguese]{babel}
\usepackage{amsmath,amssymb,amsthm}
\usepackage[margin=2.4cm]{geometry}
\usepackage{booktabs}\usepackage{xcolor}
\usepackage[hidelinks]{hyperref}

\definecolor{cinza}{gray}{0.35}
\newcommand{\code}[1]{\texttt{\small #1}}
\newcommand{\medido}[1]{\par\medskip\noindent{\small\color{cinza}\textbf{Medido.} #1}\par\medskip}
\newtheorem{teorema}{Teorema}
\newtheorem{definicao}[teorema]{Definição}
\renewcommand{\proofname}{Demonstração}

\title{\textbf{A estrela e o buraco negro}\\[2pt]
  {\large dois sinais da mesma quantidade, e o zero entre eles}}
\author{Aarão Melo Lopes}
\date{7 de agosto de 2026}

\begin{document}
\maketitle

\begin{abstract}
\noindent Deriva-se, em por-unidade e sem uma única constante física, que \textbf{a estrela é a
borda} do trial e \textbf{o buraco negro é o cristal} --- não como analogia, mas como os dois
sinais de uma mesma quantidade, com o zero entre eles. A involução $a\mapsto-a$ troca os dois lados
e \emph{fixa} a borda: é a bidualidade, e o ponto fixo é um só. Deriva-se ainda a
\textbf{pressão de radiação} contando os eixos --- três, o trial --- donde $p=\rho/3$ e $w=1/3$,
sem introduzir constante nenhuma. E mostra-se que $w=1/3$ \emph{não} pertence à família metálica,
o que é resultado e não falha.
\end{abstract}

\section{Em p.u.\ não há constantes a pedir}

$G$, $c$, $\hbar$, $\sigma$ são \textbf{factores de escala}. Em por-unidade cada um vale $1$ por
construção, e o que sobra são \textbf{razões}. Tudo o que se segue é razão --- e é por isso que
esta derivação não precisa de medir o universo para correr.

\smallskip
\noindent O equilíbrio hidrostático diz que, num ponto, o que \emph{empurra} para fora iguala o
que \emph{puxa} para dentro. Em p.u.\ isso é uma razão, e o desvio ao equilíbrio é
\[
  a \;=\; \frac{\text{empurra}}{\text{puxa}} - 1 .
\]

\section{E o sinal de $a$ é o trial}

O catálogo já classifica os regimes \emph{pelo sinal}: \emph{cristal} encolhe, \emph{borda}
conserva, \emph{caos} cresce. É o trial $\{-1,0,+1\}$, e é nele que os dois objectos moram:

\begin{center}\small
\begin{tabular}{@{}llll@{}}
\toprule
sinal & regime & o que faz & o que é \\
\midrule
$a>0$ & caos & cresce & a expansão \\
$a=0$ & \textbf{borda} & \textbf{conserva} & \textbf{a estrela} \\
$a<0$ & cristal & encolhe & \textbf{o buraco negro} \\
\bottomrule
\end{tabular}
\end{center}

\begin{teorema}[a estrela é a borda]
Com a razão em $1$, o desvio é exactamente zero e o raio não se move --- não devagar: \emph{não se
move}.
\end{teorema}

\begin{teorema}[o buraco negro é o cristal]
Com a razão abaixo de $1$, o desvio é negativo em todos os passos, e o raio cai \emph{sempre}. Não
há volta, e a razão de não haver é que \textbf{o sinal não muda}.
\end{teorema}

\section{A bidualidade: são o par}

\begin{teorema}[a involução troca os lados e fixa a borda]
A involução $a\mapsto-a$ satisfaz $\nu\circ\nu=\mathrm{id}$, leva cristal em caos e caos em
cristal, e tem \textbf{um} ponto fixo: a borda.
\end{teorema}

\noindent \textbf{Logo estrela e buraco negro não são dois objectos que se parecem}: são dois
sinais da mesma quantidade, e a borda é o zero entre eles. É o mesmo ponto fixo do trial, e a mesma
involução de toda a teoria.

\section{A pressão de radiação, derivada}

Sem introduzir constante nenhuma: a radiação reparte-se pelos \textbf{três eixos} --- o trial
outra vez --- e a pressão é a parte de \emph{um}:
\[
  p=\frac{\rho}{3}, \qquad w=\frac{p}{\rho}=\frac13 .
\]
\emph{Não se pôs $\sigma$, nem $c$, nem $a_{\text{rad}}$: contaram-se as direcções.}

\subsection*{E $w=1/3$ não é da família metálica --- e isso é resultado}

A equação de estado do catálogo é $w=\frac{2m}{3\sqrt D}-1$ com $D=m^2+4$. Impondo $w=1/3$:
\[
  16D = 36m^2 \;\Longrightarrow\; 16(m^2+4)=36m^2 \;\Longrightarrow\; 20m^2=64,
\]
que \textbf{não tem raiz inteira}. Varridos os $m$ de $-999$ a $999$: \textbf{nenhum} dá $w=1/3$.

\smallskip
\noindent \textbf{Isto não é uma falha da derivação: é onde a radiação vive.} Ela está na
\emph{fronteira}, e o catálogo já o dizia ao pôr o ponto fixo em $w=-1$. A família metálica
descreve os corpos; a radiação pura é o limite deles.

\section{Não há quarto regime}

O sinal toma três valores e não mais. Varridos $2001$ desvios: $1000$ cristal, $1$ borda, $1000$
caos, \textbf{zero} fora. \emph{O trial não tem onde pôr um quarto} --- e é por isso que estrela e
buraco negro \textbf{esgotam} o par, com a borda entre eles.

\medido{a razão em equilíbrio dá desvio $0$ e o raio fica em $1{,}000$ ao fim de $100$ passos; com
razão $0{,}900$ o raio cai em todos os $60$ passos e \emph{nunca} cresce; $\nu\circ\nu=\mathrm{id}$
em $1001$ desvios com \emph{um} ponto fixo; $p=\rho/3$ e $w=333/1000$ por contagem de eixos;
nenhum $m$ inteiro em $[-999,999]$ dá $w=1/3$; e o \textbf{controlo}: $2001$ desvios em três
regimes e nenhum fora (\code{tests/estrela.c}).}

\vfill
{\footnotesize\color{cinza}\noindent
Teoria, catálogo e medidores em
\href{https://goldenkingdom.patriatechnology.com}{goldenkingdom.patriatechnology.com}.
CC BY 4.0 (textos) $\cdot$ MIT (código).\par}

\end{document}
